\subsection{Earth and Climate Sciences}

Earth systems exhibit both long-term exponential trends and periodic forcing from seasonal and solar cycles.

\paragraph{Atmospheric $\text{CO}_2$ (Real-World):} We tested the engine against the iconic Mauna Loa Keeling Curve. The engine identified that a simple periodic sine wave was insufficient to capture the variance. It subsequently introduced an exponential anthropogenic acceleration overlaying the seasonal oscillation, decoupling human carbon emissions from the natural biospheric breathing cycle.
\paragraph{Nile River Volume (Real-World):} A notable result occurred when analyzing the historical minimum water levels of the Nile River from 622 AD to 1928 AD. The raw data contains a large, persistent reduction in variance and mean level at the end of the time series. The engine derived a sigmoidal step-change in the target mean-reverting level, identifying a regime shift centered around the year 1899. While 1899 is a known statistical changepoint in the data, the physical completion of the Aswan Low Dam did not occur until 1902. This discrepancy illustrates semantic overinterpretation: the model can recover a useful mathematical structure while attaching an unsupported historical explanation. Domain experts are therefore needed to assess the physical or historical meaning of proposed structures.
\paragraph{Solar Sunspots (Real-World):} Sunspot counts exhibit cyclic behavior, but they are highly asymmetric (fast rise, slow decay). The engine returned the $\sim 11$-year solar cycle frequency ($\omega \approx 0.57$) and introduced a quadratic non-linear term ($+c(S-S_0)^2$) to model the cycle's asymmetric shape, reducing the recorded MSE.
\paragraph{Climate and Agriculture:} The engine modeled seasonal crop-yield dependencies with multivariate inputs, generating structures that thresholded growth based on interacting temperature and rainfall limits.
